\bigskip

\begin{example}
    Define
    \[
    S \coloneqq \setbuilder*{p \in \Q^{+}}{p^{2} < 2}
    \]
    \begin{itemize}
        \item Let $p \in S$ be given. Define
        \[
        q \coloneqq p - \frac{p^{2} - 2}{p + 2} = \frac{2p + 2}{p + 2} > p
        \]
        Then
        \[
        \begin{array}{c}
            q^{2} - 2 = \frac{2(p^{2} - 2)}{(p + 2)^{2}} < 0 \\
            \Downarrow \\
            q \in S
        \end{array}
        \]
        This implies that $p \in S$ is not an upper bound of $S$.
        \item Let $p \in \Q^{+} \setminus S$ be given.
        \[
        \nexists q \in \Q^{+}, q^{2} = 2
        \]
        Therefore,
        \[
        \begin{array}{c}
            p^{2} \geq 2 \\
            \Downarrow \\
            p^{2} > 2 \\
            \Downarrow \\
            \forall x \in S, x^{2} < 2 < p^{2} \\
            \Downarrow \\
            \forall x \in S, x < p \\
            \Downarrow \\
            p \text{ is an upper bound of } S
        \end{array}
        \]
        Define
        \[
        q \coloneqq p - \frac{p^{2} - 2}{p + 2} = \frac{2p + 2}{p + 2} < p
        \]
        Then
        \[
        \begin{array}{c}
            q^{2} - 2 = \frac{2(p^{2} - 2)}{(p + 2)^{2}} > 0 \\
            \Downarrow \\
            q^{2} > 2 \\
            \Downarrow \\
            \forall x \in S, x^{2} < 2 < q^{2} \\
            \Downarrow \\
            \forall x \in S, x < q \\
            \Downarrow \\
            q \text{ is an upper bound of } S
        \end{array}
        \]
        This implies that $p$ is not a least upper bound of $S$.
    \end{itemize}
    To sum up,
    \[
    \nexists \alpha \in \Q^{+}, \alpha = \sup S
    \]
\end{example}

\bigskip

\begin{definition} \label{an:def:least-upper-bound-property}
    Let $S$ be an ordered set and $E \subseteq S$. If
    \begin{itemize}
        \item $E \neq \vnot$
        \item $E$ is bounded above
    \end{itemize}
    implies that
    \[
    \exists \alpha \in S, \alpha = \sup E
    \]
    then $S$ is said to have the \textbf{least upper bound property}.
\end{definition}

\bigskip

\begin{theorem} \label{an:thm:greatest-lower-bound-property}
    Let $S$ be an ordered set and $E \subseteq S$. Then
    \begin{itemize}
        \item $S$ has the least upper bound property
        \item $E \neq \vnot$
        \item $E$ is bounded below
    \end{itemize}
    implies that
    \[
    \exists \alpha \in S, \alpha = \inf E
    \]
\end{theorem}

\begin{proof}
    Assume that
    \begin{itemize}
        \item $S$ has the \emph{least upper bound property}
        \item $E \neq \vnot$
        \item $E$ is bounded below
    \end{itemize}
    Define
    \[
    L \coloneqq \setbuilder*{\lambda \in S}{\forall x \in E, \lambda \leq x} \subseteq S
    \]
    Then
    \[
    \begin{array}{c}
        E \text{ is bounded below} \\
        \Downarrow \\
        L \neq \vnot
    \end{array}
    \]
    Also,
    \[
    \begin{array}{c}
        E \neq \vnot \\
        \Downarrow \\
        \exists \beta \in E \\
        \Downarrow \\
        \forall \lambda \in L, \lambda \leq \beta \\
        \Downarrow \\
        L \text{ is bounded above}
    \end{array}
    \]
    By the \emph{least upper bound property},
    \[
    \exists \alpha \in S, \alpha = \sup L
    \]
    We will show that
    \[
    \alpha = \inf E
    \]
    Assume that $\alpha \notin L$. Then
    \[
    \begin{array}{c}
        \alpha \notin L \\
        \Downarrow \\
        \exists \beta \in E, \beta < \alpha \\
    \end{array}
    \]
    Since $\alpha = \sup L$,
    \[
    \begin{array}{c}
        \alpha = \sup L \\
        \Downarrow \\
        \beta \text{ is not an upper bound of } L \\
        \Downarrow \\
        \exists \gamma \in L, \beta < \gamma \leq \alpha
    \end{array}
    \]
    Then
    \[
    \begin{array}{c}
        \gamma \in L \\
        \Downarrow \\
        \forall x \in E, \gamma \leq x \\
        \Downarrow \\
        \gamma \leq \beta \\
        \Downarrow \\
        \beta < \gamma \leq \beta \\
        \Downarrow \\
        \bot
    \end{array}
    \]
    Therefore,
    \[
    \begin{array}{c}
        \alpha \in L \\
        \Downarrow \\
        \alpha \text{ is a lower bound of } E \\
    \end{array}
    \]
    Since $\alpha = \sup L$,
    \[
    \forall \delta \in S, \delta > \alpha \implies \delta \notin L \\
    \]
    To sum up,
    \begin{itemize}
        \item $\alpha$ is a lower bound of $E$
        \item $\forall \delta \in S, \delta > \alpha \implies \delta$ is not a lower bound of $E$
    \end{itemize}
    Therefore,
    \[
    \alpha = \inf E
    \]
\end{proof}


